% !TEX program = pdflatex
% Компиляция: pdflatex --shell-escape report.tex  (дважды — для оглавления)
\documentclass[a4paper, final]{article}

\usepackage{cmap}
\usepackage[russian]{babel}
\usepackage[utf8]{inputenc}
\usepackage[T1,T2A]{fontenc}
\usepackage[14pt]{extsizes}
\usepackage[left=20mm, top=20mm, right=20mm, bottom=20mm, footskip=10mm]{geometry}
\usepackage{ragged2e}
\usepackage{indentfirst}
\usepackage{listings}
\usepackage{minted}
\usepackage{hyperref}
\usepackage{graphicx}
\usepackage{caption}
\usepackage{subcaption}
\usepackage{amssymb}
\usepackage{amsmath}
\usepackage{setspace}
\usepackage{relsize}
\usepackage{enumitem}
\usepackage{longtable}
\usepackage{array}

\hypersetup{colorlinks=true,linkcolor=black,urlcolor=blue}

\newenvironment{code}{\captionsetup{type=listing}}{}

\lstset{
   language          = C,
   basicstyle        = \small\ttfamily,
   keywordstyle      = \bfseries,
   numbers           = left,
   numberstyle       = \tiny,
   numbersep         = 5pt,
   frame             = single,
   tabsize           = 2,
   inputencoding     = utf8,
   columns           = fullflexible,
   showstringspaces  = false,
   extendedchars     = false,
   breaklines        = true
}

\linespread{1.0}

% =====================================================================
\begin{document}

\thispagestyle{empty}
\begin{center}
    \hfill \break
    \hfill \break
    МИНИСТЕРСТВО НАУКИ И ВЫСШЕГО ОБРАЗОВАНИЯ РОССИЙСКОЙ ФЕДЕРАЦИИ\\[5 pt]
    Федеральное государственное автономное образовательное учреждение высшего образования
    «Санкт-Петербургский политехнический университет Петра Великого»\\[20 pt]
    Институт компьютерных наук и кибербезопасности\\ [10pt]
    Высшая школа технологий искусственного интеллекта\\[10pt]
    Направление: 02.03.01 Математика и компьютерные науки\\[80 pt]

    {\large Функциональное программирование\\ [7.5 pt]}
    {\LARGE\bf Отчет о выполнении лабораторных работ 1-5} \\
\end{center}

\vspace{60 pt}

\begin{tabular*}{460pt}{@{\extracolsep{\fill}} l r}
     Студент,\tabularnewline группы 5130201/40003 &  \line(1,0){100} \quad Лебедевич К. О.\\
     & \\
     \tabularnewline Преподаватель, \\ к.т.н. & \line(1,0){100} \quad Моторин Д.Е. \\
     & \\
     & \\
     & \\
     & \guillemotleft \line(1,0){30} \guillemotright \line(1,0){100} \, 20\,\line(1,0){15} \,г.
\end{tabular*} \\

\vspace{5 pt}
\begin{center}
    Санкт-Петербург, 2026
\end{center}

\newpage

% =========================== Оглавление ==============================
\tableofcontents
\newpage

% =========================== Введение ================================
\section{Введение}

В данном отчёте описывается процесс выполнения лабораторных работ 1--5 по
дисциплине «Функциональное программирование». Все работы выполнены на языке
Haskell. Для каждой работы приводится перечень коммитов с описанием внесённых
изменений, рассмотренные при защите вопросы и их решения. Полные листинги
исходного кода вынесены в Приложение~А.

Исходный код размещён в репозитории на GitHub с именем \texttt{hask\_pr01}.
Каждая работа находится в отдельном каталоге с именем \texttt{practice\_} и
номером работы (например, \texttt{practice\_01}).

% =================== Временная диаграмма =============================
\section{Временная диаграмма выполнения заданий}

\subsection{Пояснения к диаграмме}

В таблице~\ref{tab:timeline} приведены этапы выполнения лабораторных работ с
указанием диапазона дат активной работы (по датам коммитов). Работа над
заданиями велась последовательно с марта по май 2026 года.

\subsection{Временная диаграмма}

\begin{table}[h]
\centering
\caption{Этапы выполнения лабораторных работ}
\label{tab:timeline}
\begin{tabular}{|c|p{8cm}|c|}
\hline
\textbf{№} & \textbf{Лабораторная работа} & \textbf{Период выполнения} \\
\hline
1 & Стандартные функции, функции высшего порядка, свёртки, типы для выпекания тортов
  & 04.03 -- 27.03 \\
\hline
2 & Классы типов (\texttt{Show}, \texttt{Read}, \texttt{Eq}, \texttt{Ord},
    \texttt{Enum}, \texttt{Bounded}); собственные \texttt{MyMaybe},
    \texttt{MyEither}, \texttt{MyTree}; парсеры
  & 09.04 -- 24.04 \\
\hline
3 & Глитч-эффекты на изображениях; поиск пути в лабиринте на монаде RWS
  & 06.05 -- 08.05 \\
\hline
4 & Собственный трансформер монад \texttt{MyStateT}; текстовая игра «лабиринт»
    на стеке трансформеров
  & 14.05 -- 15.05 \\
\hline
5 & Тестирование функций с помощью QuickCheck
  & 21.05 -- 22.05 \\
\hline
\end{tabular}
\end{table}

% ============== Структура и содержание репозитория ===================
\section{Структура и содержание репозитория}

\subsection{Общая структура}

Репозиторий \texttt{hask\_pr01} организован по каталогам с именем
\texttt{practice\_} и номером работы. В каждом каталоге содержатся только
файлы и/или stack-проекты, относящиеся к соответствующему заданию.

\subsection{Содержание разделов}

\begin{enumerate}[leftmargin=*]
  \item \texttt{practice\_01} — файлы \texttt{pr01\_1.hs} и \texttt{pr01\_2.hs}
        первой работы.
  \item \texttt{practice\_02/part1/MyEvolProject} — проект первой части
        второй работы (тип \texttt{MyEvolution}).
  \item \texttt{practice\_02/part2/myParser} — stack-проект второй и третьей
        частей второй работы (модули \texttt{MyTypes} и \texttt{MyParsers}).
  \item \texttt{practice\_03/myGlithes} — первая часть третьей работы
        (глитч-эффекты).
  \item \texttt{practice\_03/myMaze} — вторая часть третьей работы
        (RWS-решатель лабиринта).
  \item \texttt{practice\_04/myMazeGame} — четвёртая работа (игра на
        трансформерах монад).
  \item \texttt{practice\_05/myFunctionsAndTests} — пятая работа
        (QuickCheck-тесты).
\end{enumerate}

% ===================================================================
% =========================== РАБОТА 1 ==============================
% ===================================================================
\section{Лабораторная работа №1}

В первой работе реализованы стандартные функции работы со списками, функции
высшего порядка, свёртки, а также система типов для процесса выпекания тортов
с контролем ингредиентов.

\subsection{Коммиты и выполняемая работа}

\subsubsection{Коммит 1. Инициализация репозитория (04.03.2026, 16:41)}

Создан репозиторий на GitHub. Добавлены первоначальные файлы первой части
задания.

\subsubsection{Коммит 2. Добавление myHead и myTail (04.03.2026, 17:05)}

Реализованы первые функции: \texttt{myHead} и \texttt{myTail} — отделение
головы и хвоста списка.

\subsubsection{Коммит 3. Стандартные функции и функции высшего порядка
               (06.03.2026, 03:59)}

Реализованы и доработаны функции первой части (\texttt{pr01\_1.hs}):
\begin{enumerate}[label=\alph*),leftmargin=*]
  \item \texttt{myHead}, \texttt{myTail} — отделение головы и хвоста;
  \item \texttt{myZip}, \texttt{myZip3} — попарное и потройное объединение;
  \item \texttt{myUnzip} — разделение списка пар на пару списков;
  \item \texttt{myFST}, \texttt{mySND}, \texttt{myTHRD} — проекции кортежа;
  \item функции высшего порядка: \texttt{myMap}, \texttt{myZipWith},
        \texttt{myZipWith3}, \texttt{myAll}, \texttt{myAny},
        \texttt{myComposition}.
\end{enumerate}

\subsubsection{Коммит 4. myZipSave и myUnzipSave (12.03.2026, 21:19)}

Во второй части (\texttt{pr01\_2.hs}) реализованы:
\begin{enumerate}[label=\alph*),leftmargin=*]
  \item \texttt{myZipSave :: [a] -> [b] -> ([(a, b)], [a], [b])} — объединение
        списков с сохранением остатка более длинного;
  \item \texttt{myUnzipSave :: ([(a, b)], [a], [b]) -> ([a], [b])} — обратная
        операция с восстановлением остатка.
\end{enumerate}

\subsubsection{Коммит 5. Свёртки и обработка Maybe (13.03.2026, 04:46)}

Реализованы функции на основе свёрток: \texttt{myFoldr}, \texttt{myFoldl},
\texttt{myFoldl1}, \texttt{myFoldr1}, \texttt{myReverse} (через свёртку),
\texttt{myTakeWhile} и \texttt{mySpan} (через \texttt{myFoldr}), а также
\texttt{myMaybe} — обработка возможно отсутствующего значения.

\subsubsection{Коммит 6. myMap для MyList и myUnfoldr (20.03.2026, 01:23)}

Определён рекурсивный тип \texttt{MyList a = MyEmpty | MyCons a (MyList a)};
для него реализована \texttt{myMap}. Добавлена развёртка
\texttt{myUnfoldr :: (b -> Maybe (a, b)) -> b -> [a]}.

\subsubsection{Коммит 7. Типы для выпекания тортов (20.03.2026, 02:28)}

Определены типы \texttt{Ingredient}, \texttt{FillingMix}, \texttt{Dough},
\texttt{CakeDoughType}, \texttt{Cake}, \texttt{Action}. Каждый конструктор
ингредиента хранит количество (\texttt{Int}), что позволяет контролировать
объём расхода.

\subsubsection{Коммит 8. Доработка контроля ингредиентов (26.03.2026, 15:02)}

Уточнены конструкторы ингредиентов (добавлен параметр \texttt{Int}); функция
\texttt{makeCakeMix} исправлена. Функция \texttt{cakeDough} откорректирована,
вспомогательная функция проверки количества удалена за счёт переноса проверок
в охранные выражения.

\subsubsection{Коммит 9. Примеры выпекания (26.03.2026, 16:31)}

Реализованы \texttt{chocolateCakeDough}, \texttt{vanillaCakeDough},
\texttt{carrotCakeDough} и \texttt{chocolateCake}, \texttt{vanillaCake},
\texttt{carrotCake}. Добавлены примеры удачного приготовления трёх видов
тортов и примеры ошибочных ситуаций: недостаток ингредиентов, неверные
ингредиенты, несоответствие смеси и теста.

\subsubsection{Коммит 10. Реорганизация каталогов (26.03.2026, 17:05)}

Файлы перенесены в каталоги \texttt{practice\_01} и \texttt{practice\_02}.

\subsection{Вопросы защиты и их решения}

\subsubsection{Вопрос 1 (задание дано 20.03.2026)}

Реализовать функцию типа \texttt{Maybe (a -> a) -> [Maybe a] -> [Maybe a]}.

\begin{code}
\begin{minted}[linenos, numbersep=5pt, frame=lines, framesep=2mm,
               fontsize=\small, breaklines=true]{haskell}
func :: Maybe (a -> a) -> [Maybe a] -> [Maybe a]
func Nothing _ = []
func _ [] = []
func (Just f) (xs) = map g xs where
    g (Just x) = Just (f x)
    g Nothing = Nothing
\end{minted}
\caption*{Листинг 1: func — применение Maybe-функции к списку Maybe-значений}
\end{code}

\newpage
\subsubsection{Вопрос 2 (задание дано 27.03.2026)}

Реализовать функцию типа
\texttt{Either String (a -> b -> c) -> Either String [(a,b)] -> Either String [c]}.

\begin{code}
\begin{minted}[linenos, numbersep=5pt, frame=lines, framesep=2mm,
               fontsize=\small, breaklines=true]{haskell}
func0 :: Either String (a -> b -> c) -> Either String [(a,b)] -> Either String [c]
func0 (Left a) (Left b) = Left (a ++ b)
func0 (Left a) _ = Left a
func0 _ (Left b) = Left b
func0 (Right f) (Right xs) = Right (helper f xs) where
    helper f ((a,b): xxs) = (f a b) : helper f xxs
    helper f [] = []
\end{minted}
\caption*{Листинг 2: func0 — применение Either-функции к списку пар Either}
\end{code}

% ===================================================================
% =========================== РАБОТА 2 ==============================
% ===================================================================
\section{Лабораторная работа №2}

Вторая работа состоит из трёх частей: реализация классов типов для
перечислимого типа \texttt{MyEvolution}; собственные типы \texttt{MyMaybe},
\texttt{MyEither}, \texttt{MyTree} как представители стандартных классов
типов; парсеры выражений (собственный и на основе библиотек Parsec и
Attoparsec).

\subsection{Часть 1. Тип MyEvolution}

\subsubsection{Коммит 1. Проект MyEvolProject и тип MyEvolution (09.04.2026, 19:02)}

Создан проект \texttt{MyEvolProject}. Определён тип-сумма
\texttt{MyEvolution} с десятью конструкторами — этапами эволюции (от
\texttt{LUCA} до \texttt{Humans}). Вручную реализованы представители классов:
\begin{enumerate}[label=\alph*),leftmargin=*]
  \item \texttt{Show} — каждый конструктор преобразуется в строку (например,
        \texttt{LUCA} $\to$ «Last Universal Common Ancestor»);
  \item \texttt{Read} — обратное преобразование строки в конструктор;
  \item \texttt{Eq} — равенство одинаковых конструкторов;
  \item \texttt{Enum} — \texttt{fromEnum}/\texttt{toEnum} для порядковых
        номеров;
  \item \texttt{Ord} — сравнение через \texttt{fromEnum};
  \item \texttt{Bounded} — \texttt{minBound = LUCA}, \texttt{maxBound = Humans}.
\end{enumerate}
Дополнительно определён тип \texttt{MyEvolution'} с теми же конструкторами,
представители классов получены автоматически через \texttt{deriving}. В
\texttt{Main.hs} выводится отсортированный список всех значений типа.

\subsection{Часть 2. Типы MyMaybe, MyEither, MyTree}

\subsubsection{Коммит 2. MyTree: Functor (10.04.2026, 00:04)}

Определён тип \texttt{MyTree a = Leaf a | Node a (MyTree a) (MyTree a)},
хранящий значения и в листьях, и в узлах. Реализован представитель
\texttt{Functor}.

\subsubsection{Коммит 3. MyTree: Foldable (15.04.2026, 02:42)}

Реализован представитель \texttt{Foldable} для \texttt{MyTree} (\texttt{foldr},
\texttt{foldMap}, \texttt{length}). Добавлены примеры вызовов в комментариях.

\subsubsection{Коммит 4. MyEither (16.04.2026, 03:05 )}

Определён тип \texttt{MyEither e a = MyLeft e | MyRight a}. Реализованы
представители \texttt{Functor}, \texttt{Foldable}, \texttt{Applicative} и
\texttt{Semigroup}. Добавлены примеры вызовов всех реализованных функций
(\texttt{<>}, \texttt{stimes}, \texttt{sconcat}, \texttt{(*>)},
\texttt{(<*)}, \texttt{liftA2}, \texttt{(<\$)}, \texttt{fold} и другие).

\subsubsection{Коммит 5. MyTree: Applicative (15.04.2026, 04:13)}

Реализован представитель \texttt{Applicative} для \texttt{MyTree}
(\texttt{pure}, \texttt{<*>}) и примеры вызовов \texttt{<*>},
\texttt{liftA2}, \texttt{*>}, \texttt{<*}.

\subsubsection{Коммит 6. MyMaybe (15.04.2026, 05:24)}

Определён тип \texttt{MyMaybe a = MyJust a | MyNothing}. Реализованы
представители \texttt{Functor}, \texttt{Applicative}, \texttt{Semigroup},
\texttt{Monoid} (с \texttt{mempty = MyNothing}) и \texttt{Foldable}. В
комментариях приведены примеры вызовов \texttt{fmap}, \texttt{<\$},
\texttt{pure}, \texttt{<*>}, \texttt{liftA2}, \texttt{<>}, \texttt{sconcat},
\texttt{stimes}, \texttt{mappend}, \texttt{mconcat}, \texttt{foldr},
\texttt{foldMap}, \texttt{fold}.

\subsection{Часть 3. Парсеры}

\subsubsection{Коммит 7. Собственный парсер MyParser (23.04.2026, 15:51)}

Реализован модуль \texttt{MyParsers.MyParser}. Определён тип
\texttt{newtype Parser tok a = Parser \{ runParser :: [tok] -> MyMaybe ([tok], a) \}},
использующий собственный тип \texttt{MyMaybe} вместо стандартного
\texttt{Maybe}. Реализованы представители \texttt{Functor},
\texttt{Applicative}, \texttt{Alternative}, базовые парсеры
(\texttt{satisfy}, \texttt{char}, \texttt{lower}, \texttt{digit}) и
комбинированные (\texttt{multiplication}, \texttt{lowers}, \texttt{digits},
\texttt{finalMult}, \texttt{finalPlus}, \texttt{plusOrMult}). В
\texttt{Main.hs} добавлен квалифицированный импорт всех трёх парсеров.

\subsubsection{Коммит 8. ParsecParser (23.04.2026, 20:03)}

Реализован модуль \texttt{MyParsers.ParsecParser} на основе библиотеки
Parsec: парсеры \texttt{digitP}, \texttt{multiplicationP}, \texttt{digitsP},
\texttt{finalMultP}, \texttt{finalPlusP}, \texttt{plusOrMultParsec} (с
использованием \texttt{try} для backtracking). Функция
\texttt{runParsecParser} запускает парсер и возвращает результат в виде
\texttt{MyMaybe}.

\subsubsection{Коммит 9. AttoparsecParser (23.04.2026, 21:19)}

Реализован модуль \texttt{MyParsers.AttoparsecParser} на основе библиотеки
Attoparsec: аналогичные парсеры \texttt{digitA}, \texttt{multiplicationA},
\texttt{digitsA}, \texttt{finalMultA}, \texttt{finalPlusA},
\texttt{plusOrMultAttoparsec}. Функция \texttt{runParser} преобразует
\texttt{String} в \texttt{Text} и возвращает результат в виде \texttt{MyMaybe}.
Добавлен вывод результата для Attoparsec.

\subsubsection{Коммит 10. Комментарии-промпты (24.04.2026, 11:40)}

Добавлены комментарии-промпты, использованные при генерации модулей Parsec и
Attoparsec.

\subsection{Вопросы защиты и их решения}

\subsubsection{Вопрос 1 (задание дано 24.04.2026)}

Составить выражение с \texttt{<\$>}, \texttt{<*>}, \texttt{<*>} типа
\texttt{Either a (String, Float)}.

Решение:

\begin{code}
\begin{minted}[linenos, numbersep=5pt, frame=lines, framesep=2mm,
               fontsize=\small, breaklines=true]{haskell}
(,) <$> (Right "abc") <*> (Right (+1.0) <*> (Right (4.0 :: Float)))
\end{minted}
\caption*{Листинг 3: Either a (String, Float) через аппликативный функтор}
\end{code}

Первый \texttt{<*>} применяет функцию \texttt{(+1.0)} к значению \texttt{4.0},
второй — конструктор пары \texttt{(,)} к \texttt{"abc"} и полученному
\texttt{5.0}. Результат: \texttt{Right ("abc", 5.0)}.

\subsubsection{Вопрос 2 (задание дано 08.05.2026)}

Составить выражение с \texttt{<\$>}, \texttt{<*>}, \texttt{<*>} типа
\texttt{Either a (Int, String)}.

Решение:

\begin{code}
\begin{minted}[linenos, numbersep=5pt, frame=lines, framesep=2mm,
               fontsize=\small, breaklines=true]{haskell}
(,) <$> (Right (1 :: Integer)) <*> (Right (++ "1") <*> (Right ("abc" :: String)))
\end{minted}
\caption*{Листинг 4: Either a (Int, String) через аппликативный функтор}
\end{code}

Первый \texttt{<*>} применяет функцию \texttt{(++ "1")} к строке
\texttt{"abc"}, второй — конструктор пары к \texttt{1} и \texttt{"abc1"}.
Результат: \texttt{Right (1, "abc1")}.

% ===================================================================
% =========================== РАБОТА 3 ==============================
% ===================================================================
\section{Лабораторная работа №3}

Третья работа состоит из двух частей: генерация глитч-эффектов на изображениях
и поиск пути в лабиринте с использованием монады RWS.

\subsection{Часть 1. Глитч-эффекты}

\subsubsection{Коммит 1. Создание проектов (06.05.2026, 17:36)}

Добавлены файлы stack-проектов \texttt{myGlithes} и \texttt{myMaze}.

\subsubsection{Коммит 2. Функции глитч-эффектов (07.05.2026, 13:17)}

В модуле \texttt{Glitches} реализованы:
\begin{enumerate}[label=\alph*),leftmargin=*]
  \item \texttt{intToChar} — преобразование \texttt{Int} в \texttt{Char} по
        модулю 255;
  \item \texttt{intToBC} — преобразование \texttt{Int} в \texttt{ByteString}
        из одного символа;
  \item \texttt{replaceByte} — детерминированная замена байта в заданной позиции;
  \item \texttt{randomReplaceByte} — замена случайного байта на случайное
        значение;
  \item \texttt{sortSection} — сортировка участка \texttt{ByteString} по
        заданному смещению и размеру;
  \item \texttt{randomSortSection} — сортировка случайного участка
        фиксированного размера (25 байт).
\end{enumerate}
В \texttt{Main.hs} имя файла берётся из аргументов командной строки, к
изображению последовательно применяется цепочка глитч-функций (через
\texttt{foldM}), каждая версия сохраняется в отдельный файл.

\subsection{Часть 2. Лабиринт на монаде RWS}

\subsubsection{Коммит 3. Типы данных лабиринта (07.05.2026, 18:50)}

В модуле \texttt{Maze} определены типы:
\texttt{RoomType = Start String | Finish String | Node String},
\texttt{Room = Room RoomType [String]} (тип комнаты и список соседей),
\texttt{Maze = [Room]}. Реализованы функции \texttt{getName},
\texttt{getNeighbours}, \texttt{isFinish}, \texttt{findFinish},
\texttt{findStart}.

\subsubsection{Коммит 4. Поиск пути (08.05.2026, 00:24)}

Добавлены функции \texttt{findRoom}, \texttt{isInPath},
\texttt{neighboursCount}, \texttt{getUnvisited}, а также основной
решатель \texttt{solveMaze} на монаде RWS: окружение — лабиринт, состояние —
пройденный путь, лог — журнал перемещений. Реализован вспомогательный обход
\texttt{searchPath} (поиск в глубину) с логированием и обработкой ошибок
(отсутствие старта/финиша, тупик). По требованию задания do-нотация не
используется.

\subsubsection{Коммит 5. Тестирование и очистка (08.05.2026, 02:33)}

В \texttt{Main.hs} добавлены тестовые лабиринты \texttt{schoolMaze} и
\texttt{castleMaze} с выводом результата, пути и логов. Функции
\texttt{findStart}/\texttt{findFinish} переработаны для устранения
предупреждений компилятора; из исходников удалена кириллица.

\subsection{Вопросы защиты и их решения}

\subsubsection{Вопрос 1 (задание дано 08.05.2026)}

Составить выражение вида \texttt{(\_ >>=\ \_ >>=\ \_)} с монадой \texttt{State}.

Решение:

\begin{code}
\begin{minted}[linenos, numbersep=5pt, frame=lines, framesep=2mm,
               fontsize=\small, breaklines=true]{haskell}
f1 = state (\x -> ("x", "log"))
f2 = (\x -> state (\x -> ("x", "log")))
runState (f1 >>= f2 >>= f2) "abc"
-- ("x","log")
\end{minted}
\caption*{Листинг 5: Цепочка >>= для монады State}
\end{code}


\section{Лабораторная работа №4}

В четвёртой работе реализован собственный трансформер монад \texttt{MyStateT}
и текстовая игра «прохождение лабиринта» на стеке из трёх трансформеров.

\subsection{Коммиты и выполняемая работа}

\subsubsection{Коммит 1. Трансформер MyStateT (14.05.2026, 22:35)}

Определён тип
\texttt{newtype MyStateT s m a = MyStateT \{ runMyStateT :: s -> m (a, s) \}}.
Реализована функция \texttt{myState} (поднятие чистой функции преобразования
состояния) и представители классов:
\begin{enumerate}[label=\alph*),leftmargin=*]
  \item \texttt{Functor} (при условии \texttt{Functor m});
  \item \texttt{Applicative} (при условии \texttt{Monad m});
  \item \texttt{Monad} (при условии \texttt{Monad m});
  \item \texttt{MonadTrans} с методом \texttt{lift}.
\end{enumerate}

\subsubsection{Коммит 2. Загрузка лабиринта (14.05.2026, 23:39)}

В модуль \texttt{Maze} перенесены типы и вспомогательные функции из третьей
работы. Добавлены \texttt{parseRoom :: String -> Room} и
\texttt{loadMaze :: FilePath -> IO Maze} для чтения карты лабиринта из
текстового файла. Старт и финиш распознаются по именам «старт» и «финиш».

\subsubsection{Коммит 3. Игровая логика (15.05.2026, 01:32)}

Определён стек трансформеров
\texttt{type Game a = ReaderT Maze (WriterT [String] (MyStateT String IO)) a}:
\texttt{ReaderT} — доступ к лабиринту, \texttt{WriterT} — логирование,
\texttt{MyStateT} — текущая комната, \texttt{IO} — ввод-вывод. Реализованы
функции \texttt{logPrint}, \texttt{getCRoom}, \texttt{setCRoom},
\texttt{runGame} (с явной терминацией трансформеров), основной игровой цикл
\texttt{gameProc} и точка входа \texttt{playMaze}. Игрок перемещается по
комнатам, может завершить игру вводом «q»; перемещения логируются и выводятся
в конце.

\subsubsection{Коммит 4. Тестовый лабиринт (15.05.2026, 01:33)}

Добавлен файл \texttt{mazes/uni.maze} с тестовой картой лабиринта.

\subsubsection{Коммит 5. Комментарии к runGame (15.05.2026, 17:56)}

Добавлены подробные комментарии к функции \texttt{runGame}, поясняющие
пошаговое развёртывание (терминацию) стека трансформеров.

\subsection{Вопросы защиты и ответы на них}

\subsubsection{Вопрос 1}

Как работают типы в строке
\texttt{runMyStateT (runWriterT (runReaderT game maze)) startRoom}?

Решение: применение функций снимает трансформеры по очереди:
\begin{enumerate}[label=\arabic*),leftmargin=*]
  \item \texttt{runReaderT game maze} даёт
        \texttt{WriterT [String] (MyStateT String IO) a};
  \item \texttt{runWriterT (\ldots)} даёт
        \texttt{MyStateT String IO (a, [String])};
  \item \texttt{runMyStateT (\ldots) startRoom} даёт
        \texttt{IO ((a, [String]), String)}.
\end{enumerate}


\subsubsection{Вопрос 2}

Что такое терминатор?

Решение: терминатор — это самая нижняя монада в стеке трансформеров, то есть
базовая монада, поверх которой надстраиваются все остальные слои. В данном
проекте терминатором является \texttt{IO}:

\subsubsection{Вопрос 3}

Что такое монадный трансформер?

Решение: монадный трансформер — это тип вида \texttt{T m a}, параметризованный
другой монадой \texttt{m}, который добавляет к \texttt{m} новый эффект, сохраняя возможность
использовать эффекты \texttt{m}. Трансформер реализует класс
\texttt{MonadTrans} с методом \texttt{lift :: m a -> T m a}, позволяющим
«поднимать» вычисления внутренней монады в расширенный стек.

\subsubsection{Вопрос 4}

Где запускается монада \texttt{IO}?

Решение: монада \texttt{IO} запускается через оператор \texttt{(<-)} в
do-нотации.


% ===================================================================
% =========================== РАБОТА 5 ==============================
% ===================================================================
\section{Лабораторная работа №5}

В пятой работе реализованы вспомогательные функции и набор свойств для их
тестирования с помощью библиотеки QuickCheck.

\subsection{Коммиты и выполняемая работа}

\subsubsection{Коммит 1. Создание проекта (21.05.2026, 22:50)}

Создан stack-проект \texttt{myFunctionsAndTests}, в \texttt{package.yaml}
добавлена зависимость \texttt{QuickCheck} для тестов.

\subsubsection{Коммит 2. Тестируемые функции (21.05.2026, 23:27)}

В модуле \texttt{Lib} реализованы:
\begin{enumerate}[label=\alph*),leftmargin=*]
  \item \texttt{normMod :: Int -> Int -> Int} — приведение числа к остатку по
        модулю через повторное вычитание/сложение (вспомогательная функция);
  \item \texttt{addMod :: Int -> Int -> Int -> Int} — сложение по модулю;
  \item \texttt{splitTokens :: String -> [String]} — разбиение строки на
        чередующиеся блоки слов и пробелов;
  \item \texttt{clean :: [String] -> [String]} — удаление пустых токенов;
  \item \texttt{reverseWords :: String -> String} — изменение порядка слов
        на обратный (реализована как \texttt{concat . reverse . splitTokens}).
\end{enumerate}

\subsubsection{Коммит 3. Генераторы и свойства (22.05.2026, 01:24)}

В \texttt{test/Spec.hs} реализованы генераторы \texttt{genMod},
\texttt{genWord}, \texttt{genSpaces}, \texttt{genPhrase}, \texttt{genText}
и свойства:
\begin{enumerate}[label=\alph*),leftmargin=*]
  \item \texttt{prop\_addMod\_mod} — согласованность результата с
        \texttt{mod}: \texttt{addMod x y m \`mod\` m == (x + y) \`mod\` m};
  \item \texttt{prop\_addMod\_neutral} — нейтральный элемент:
        \texttt{addMod x 0 m == x \`mod\` m};
  \item \texttt{prop\_addMod\_comm} — коммутативность:
        \texttt{addMod x y m == addMod y x m};
  \item \texttt{prop\_rev\_empty} — переворот пустой строки даёт пустую
        строку;
  \item \texttt{prop\_rev\_single} — строка из одного слова не меняется;
  \item \texttt{prop\_rev\_order} — порядок слов меняется на обратный;
  \item \texttt{prop\_rev\_involution} — двойное применение возвращает
        исходную строку.
\end{enumerate}

\subsubsection{Коммит 4. verbose и withMaxSuccess (22.05.2026, 01:29)}

В \texttt{main} тесты обёрнуты в \texttt{verbose} и \texttt{withMaxSuccess 5}
для подробного вывода ограниченного числа проверок.


\newpage
\appendix
\section{Приложение А. Листинги кода}

% Полные листинги исходного кода будут добавлены отдельно.

\end{document}
